% Chapter — Rigid-body attitude dynamics (FR-1). Follows the mandatory
% FR-29 six-section template: purpose; assumptions and domain of validity;
% derivation; discretization and implementation; validation evidence;
% references. The equation labels eq:rigidbody:qdot, eq:rigidbody:euler,
% eq:rigidbody:omegadot, and eq:rigidbody:normdrift are echoed verbatim in
% comments in cpp/src/models/rigidbody.cpp and cpp/tests/test_rigidbody.cpp
% (FR-29 equation-label-to-code traceability); renaming them breaks that
% trace.
\chapter{Rigid-Body Attitude Dynamics}
\label{ch:rigidbody}

\section{Purpose}
\label{sec:rigidbody:purpose}

This chapter documents the rotational core of the FR-1 six-degree-of-freedom
vehicle model: the quaternion attitude kinematics and Euler's equation for a
rigid body with time-varying inertia, including the $\dot{\mat{I}}\vv{\omega}$
term. The module supplies the attitude piece of the vehicle equations of
motion as a pure right-hand-side map: given the attitude quaternion
$\quat{q}_{i2b}$, the body rate $\vv{\omega}^{B}$, the inertia tensor
$\mat{I}$, its time derivative $\dot{\mat{I}}$, and the total body torque
$\vv{\tau}^{B}$, it returns $(\dot{\quat{q}}, \dot{\vv{\omega}})$ for
embedding in a larger state vector. The Phase~4 vehicle layer supplies
$\mat{I}$, $\dot{\mat{I}}$, and $\vv{\tau}^{B}$ from mass properties,
actuators, and environmental torques (Chapter~\ref{ch:gravgrad}); this module
depends on none of those structures. It is implemented in
\texttt{cpp/src/models/rigidbody.cpp} with its interface in
\texttt{cpp/include/star/models/rigidbody.hpp}. The chapter also derives the
quaternion norm-drift bound that the FR-1 post-step normalization policy
rests on.

\section{Assumptions and domain of validity}
\label{sec:rigidbody:domain}

Assumptions:
\begin{enumerate}
  \item \textbf{Quasi-rigid body.} The vehicle is rigid at every instant;
    mass variation enters only through the externally supplied pair
    $\bigl(\mat{I}(t), \dot{\mat{I}}(t)\bigr)$, both expressed in body axes
    about the \emph{instantaneous} center of mass, with $\mat{I}$ symmetric
    positive definite and the pair mutually consistent (the Phase~4
    mass-properties model owns that consistency). Momentum carried by
    internal relative motion of the depleting mass (propellant swirl, jet
    damping) is excluded: the $\dot{\mat{I}}\vv{\omega}$ form is the
    standard quasi-rigid depletion approximation.
  \item \textbf{Torque completeness.} $\vv{\tau}^{B}$ is the total external
    plus actuator torque about the instantaneous center of mass, resolved in
    body axes. The module does not distinguish torque sources.
  \item \textbf{Attitude convention.} $\quat{q}_{i2b}$ is the Hamilton,
    scalar-first, inertial-to-body frame-transformation quaternion of
    Section~\ref{sec:notation:quaternion} (decision D-7);
    $\vv{\omega}^{B}$ is the angular velocity of $B$ with respect to $I$,
    resolved in $B$.
  \item \textbf{Norm handling.} The kinematics of
    equation~\eqref{eq:rigidbody:qdot} is degree-one homogeneous in
    $\quat{q}$ and preserves $\norm{\quat{q}}$ exactly
    (equation~\eqref{eq:rigidbody:qnorm}), so the right-hand side is valid
    for any nonzero quaternion. Consumers of the attitude (the DCM of
    equation~\eqref{eq:notation:quat2dcm}) require unit norm; the post-step
    normalization policy of Section~\ref{sec:rigidbody:implementation}
    maintains it.
\end{enumerate}

Domain of validity: any state with $\mat{I}$ invertible. Out-of-domain
behavior, matching the code exactly: a singular $\mat{I}$ makes the
closed-form $3 \times 3$ inverse divide by a zero determinant and IEEE
non-finite values propagate; no right-hand-side function of this module
throws, because it runs inside the deterministic time loop and excluding
non-physical mass properties is the configuration validator's job. The one
between-step helper, post-step renormalization, throws
\texttt{std::domain\_error} on a zero or non-finite quaternion norm rather
than fabricating an attitude, mirroring
\texttt{star::rotation::quat\_normalize}: a state that degenerate means the
integration has already failed, and the deterministic response is a loud
abort.

\section{Derivation}
\label{sec:rigidbody:derivation}

\subsection{Quaternion kinematics}

Differentiate the coordinate-transformation sandwich of
equation~\eqref{eq:notation:sandwich} for a vector fixed in the inertial
frame, $\dot{\tilde{\vv{v}}}^{I} = 0$, using
$\mathrm{d}(\quat{q}^{-1})/\mathrm{d}t = -\quat{q}^{-1} \quatmul
\dot{\quat{q}} \quatmul \quat{q}^{-1}$:
\begin{equation}
  \dot{\tilde{\vv{v}}}^{B}
    = -\quat{q}^{-1} \dot{\quat{q}} \quatmul \tilde{\vv{v}}^{B}
      + \tilde{\vv{v}}^{B} \quatmul \quat{q}^{-1} \dot{\quat{q}} .
  \label{eq:rigidbody:sandwichdot}
\end{equation}
A vector fixed in $I$ appears in $B$ to rotate with
$-\vv{\omega}^{B}$: $\dot{\vv{v}}^{B} = -\vv{\omega}^{B} \times
\vv{v}^{B}$. Setting $\quat{q}^{-1} \dot{\quat{q}} =
\tfrac{1}{2}\tilde{\vv{\omega}}$ with the pure quaternion
$\tilde{\vv{\omega}} = [\,0, \vv{\omega}^{B}\,]$ makes the right side of
equation~\eqref{eq:rigidbody:sandwichdot} equal
$\tfrac{1}{2}(\tilde{\vv{v}} \quatmul \tilde{\vv{\omega}} -
\tilde{\vv{\omega}} \quatmul \tilde{\vv{v}})$, which by the Hamilton
product~\eqref{eq:notation:hamiltonproduct} is the pure quaternion with
vector part $\vv{v}^{B} \times \vv{\omega}^{B}$ --- exactly the required
rate. Hence the attitude kinematics in the project convention is
\begin{equation}
  \boxed{\;
  \dot{\quat{q}}_{i2b}
    = \tfrac{1}{2}\, \quat{q}_{i2b} \quatmul
      [\,0, \vv{\omega}^{B}\,] ,
  \;}
  \label{eq:rigidbody:qdot}
\end{equation}
the standard result (Markley and Crassidis~\cite{markley2014}, ch.~3)
transcribed to the Hamilton scalar-first $\quat{q}_{i2b}$ convention of
Chapter~\ref{ch:notation}. The norm is a first integral of
equation~\eqref{eq:rigidbody:qdot} for \emph{any} $\quat{q}$: the scalar
part of $\quat{q}^{*} \quatmul (\quat{q} \quatmul \tilde{\vv{\omega}}) =
\norm{\quat{q}}^{2}\, \tilde{\vv{\omega}}$ is zero, so
\begin{equation}
  \frac{\mathrm{d}}{\mathrm{d}t} \norm{\quat{q}}^{2}
    = 2\, \quat{q} \cdot \dot{\quat{q}} = 0 .
  \label{eq:rigidbody:qnorm}
\end{equation}
Homogeneity has a second consequence used below: if $\quat{q}(t)$ solves
equation~\eqref{eq:rigidbody:qdot} then so does $c\,\quat{q}(t)$ for any
constant $c > 0$, and both represent the same attitude history, so a small
norm error is a pure bookkeeping error --- it never feeds back into the
attitude or the rates.

\subsection{Euler's equation with time-varying inertia}

The angular momentum about the center of mass is $\vv{H}^{B} =
\mat{I}\vv{\omega}^{B}$, and Newton--Euler dynamics equates its
\emph{inertial} derivative to the applied torque. The transport theorem
between the inertial and body derivatives of any vector gives
$\dot{\vv{H}}\big|_{I} = \dot{\vv{H}}\big|_{B} + \vv{\omega} \times
\vv{H}$, so
\begin{equation}
  \mat{I}\dot{\vv{\omega}}
    + \dot{\mat{I}}\vv{\omega}
    + \vv{\omega} \times (\mat{I}\vv{\omega})
    = \vv{\tau}^{B} ,
  \qquad\text{i.e.}\qquad
  \mat{I}\dot{\vv{\omega}}
    = \vv{\tau}^{B}
      - \vv{\omega} \times (\mat{I}\vv{\omega})
      - \dot{\mat{I}}\vv{\omega} ,
  \label{eq:rigidbody:euler}
\end{equation}
(Markley and Crassidis~\cite{markley2014}, ch.~3, with the
$\dot{\mat{I}}\vv{\omega}$ term retained per FR-1). The implemented form
solves the $3 \times 3$ system explicitly:
\begin{equation}
  \boxed{\;
  \dot{\vv{\omega}}
    = \mat{I}^{-1}\bigl( \vv{\tau}^{B}
      - \vv{\omega} \times (\mat{I}\vv{\omega})
      - \dot{\mat{I}}\vv{\omega} \bigr) .
  \;}
  \label{eq:rigidbody:omegadot}
\end{equation}

For $\vv{\tau} = 0$ and $\dot{\mat{I}} = 0$ the motion has two first
integrals, the acceptance instruments of the intermediate-axis test:
\begin{align}
  \vv{H}^{I} &= \bigl(\dcm{I}{B}\bigr)^{\mathsf T} \mat{I}\vv{\omega}^{B}
    = \text{const}, \label{eq:rigidbody:H} \\
  T &= \tfrac{1}{2}\, \vv{\omega}^{\mathsf T} \mat{I} \vv{\omega}
    = \text{const}, \label{eq:rigidbody:T}
\end{align}
the first because $\dot{\vv{H}}\big|_{I} = \vv{\tau} = 0$, the second
because $\dot{T} = \vv{\omega} \cdot \mat{I}\dot{\vv{\omega}} =
\vv{\omega} \cdot (\vv{\tau} - \vv{\omega} \times (\mat{I}\vv{\omega})) =
\vv{\omega} \cdot \vv{\tau} = 0$.

\subsection{Closed-form torque-free motion of an axisymmetric body}
\label{sec:rigidbody:coning}

Let $\mat{I} = \mathrm{diag}(I_T, I_T, I_A)$ in body principal axes with
symmetry axis $\hat{\vv{e}}_3$, $\vv{\tau} = 0$, $\dot{\mat{I}} = 0$.
Equation~\eqref{eq:rigidbody:euler} componentwise gives $\dot{\omega}_3 =
0$ and $\dot{\omega}_1 = -\omega_p \omega_2$, $\dot{\omega}_2 = \omega_p
\omega_1$ with the constant relative-spin (body precession) rate
\begin{equation}
  \omega_p = \frac{I_A - I_T}{I_T}\, \omega_3 .
  \label{eq:rigidbody:omegap}
\end{equation}
The attitude solution is classical (Markley and
Crassidis~\cite{markley2014}, ch.~3); it is re-derived here in the project
convention because every sign in it is load-bearing for the golden vectors.
Since $\vv{H}^{I}$ is fixed, define an auxiliary inertial frame $N$ with
$\hat{\vv{z}}_N = \vv{H}^{I}/H$, $H = \norm{\vv{H}}$, and parametrize
$\dcm{N}{B} = \mat{R}_3(\psi)\mat{R}_1(\theta)\mat{R}_3(\phi)$ with the
3-1-3 angles of equation~\eqref{eq:rotations:e313}. Resolving $\vv{H} = H
\hat{\vv{z}}_N$ in body axes and equating to
$\mat{I}\vv{\omega}^{B}$:
\begin{equation}
  H \begin{bmatrix} \sin\theta\sin\psi \\ \sin\theta\cos\psi \\
     \cos\theta \end{bmatrix}
  = \begin{bmatrix} I_T\,\omega_1 \\ I_T\,\omega_2 \\ I_A\,\omega_3
    \end{bmatrix}
  \quad\Longrightarrow\quad
  \cos\theta = \frac{I_A\,\omega_3}{H}, \qquad
  \tan\psi = \frac{\omega_1}{\omega_2},
  \label{eq:rigidbody:coningparams}
\end{equation}
so the nutation angle $\theta$ is constant. For the chained rotation
$\quat{q}_{n2b} = \quat{q}_z(\phi) \quatmul \quat{q}_x(\theta) \quatmul
\quat{q}_z(\psi)$ with $\dot{\theta} = 0$, angular-velocity addition
(each elementary rate about its own current axis, resolved in $B$;
Markley and Crassidis~\cite{markley2014}, ch.~3) gives the 3-1-3 rate map
\begin{equation}
  \vv{\omega}^{B}
  = \dot{\phi}
    \begin{bmatrix} \sin\theta\sin\psi \\ \sin\theta\cos\psi \\
      \cos\theta \end{bmatrix}
    + \dot{\psi}
    \begin{bmatrix} 0 \\ 0 \\ 1 \end{bmatrix} .
  \label{eq:rigidbody:eulerrates}
\end{equation}
Matching equations~\eqref{eq:rigidbody:coningparams}
and~\eqref{eq:rigidbody:eulerrates} component by component yields constant
rates
\begin{equation}
  \dot{\phi} = \frac{H}{I_T},
  \qquad
  \dot{\psi} = \omega_3 - \dot{\phi}\cos\theta
             = \omega_3 \Bigl(1 - \frac{I_A}{I_T}\Bigr)
             = -\,\omega_p ,
  \label{eq:rigidbody:coningrates}
\end{equation}
consistent with the Euler-equation solution above. The full attitude
history from an arbitrary initial quaternion $\quat{q}_0$ (which fixes the
constant offset $\quat{q}_{i2n}$, with the free angle $\phi_0 = 0$
absorbing the choice of $\hat{\vv{x}}_N$) is
\begin{equation}
  \boxed{\;
  \quat{q}_{i2b}(t)
    = \quat{q}_{i2n} \quatmul
      \quat{q}_{313}\bigl(\dot{\phi}\,t,\; \theta,\;
        \psi_0 + \dot{\psi}\,t\bigr),
  \qquad
  \quat{q}_{i2n} = \quat{q}_0 \quatmul
      \quat{q}_{313}(0, \theta, \psi_0)^{-1} .
  \;}
  \label{eq:rigidbody:coning}
\end{equation}
This is the reference for the coning golden vectors
(\texttt{tests/golden/attitude/coning.toml}); the generator additionally
verifies equations~\eqref{eq:rigidbody:coning}
and~\eqref{eq:rigidbody:eulerrates} against
equations~\eqref{eq:rigidbody:qdot} and~\eqref{eq:rigidbody:euler} by
extended-precision central finite differences before writing, so the
committed reference is machine-checked to be the unique solution of the
committed initial-value problem.

\subsection{Instability of intermediate-axis spin}

For distinct principal moments $I_1 > I_2 > I_3$, linearize
equation~\eqref{eq:rigidbody:euler} (torque-free, constant $\mat{I}$)
about the intermediate-axis spin $\vv{\omega} = (0, \omega_2, 0)$:
$\dot{\omega}_1 = \frac{I_2 - I_3}{I_1}\omega_2\,\omega_3$ and
$\dot{\omega}_3 = \frac{I_1 - I_2}{I_3}\omega_1\,\omega_2$ with $\omega_2$
constant to first order, so $\ddot{\omega}_1 = \lambda^2 \omega_1$ with
the real growth rate
\begin{equation}
  \lambda = |\omega_2| \sqrt{\frac{(I_2 - I_3)(I_1 - I_2)}{I_1 I_3}} > 0 .
  \label{eq:rigidbody:lambda}
\end{equation}
A small transverse perturbation therefore grows as $e^{\lambda t}$ until
the nonlinear terms produce the finite flip (the Dzhanibekov effect). Two
testing consequences, recorded in the golden manifest: state-level
reference checkpoints are meaningful only over horizons where
$e^{\lambda t}$ has not yet amplified the checkpoint tolerance to the gate
(here $t \le \qty{15}{\second}$, $e^{\lambda t} \le 49$), and across the
flip itself only the invariants~\eqref{eq:rigidbody:H}
and~\eqref{eq:rigidbody:T} are testable --- which is exactly what Phase~4
exit criterion~4 gates.

\subsection{Quaternion norm drift under Runge--Kutta integration}
\label{sec:rigidbody:normdrift}

In the packed coordinates $\vv{y}_q = [q_w, q_x, q_y, q_z]^{\mathsf T}$,
equation~\eqref{eq:rigidbody:qdot} is the linear system $\dot{\vv{y}}_q =
\tfrac{1}{2}\mat{\Omega}(\vv{\omega})\,\vv{y}_q$ with the skew-symmetric
right-multiplication matrix
\begin{equation}
  \mat{\Omega}(\vv{\omega}) =
  \begin{bmatrix}
    0 & -\omega_x & -\omega_y & -\omega_z \\
    \omega_x & 0 & \omega_z & -\omega_y \\
    \omega_y & -\omega_z & 0 & \omega_x \\
    \omega_z & \omega_y & -\omega_x & 0
  \end{bmatrix},
  \qquad
  \mat{\Omega}^2 = -\norm{\vv{\omega}}^2 \mat{I}_4 ,
  \label{eq:rigidbody:bigomega}
\end{equation}
where the square follows from associativity: $(\quat{q} \quatmul
\tilde{\vv{\omega}}) \quatmul \tilde{\vv{\omega}} = \quat{q} \quatmul
(\tilde{\vv{\omega}} \quatmul \tilde{\vv{\omega}}) =
-\norm{\vv{\omega}}^2 \quat{q}$. For \emph{constant} $\vv{\omega}$, one
classical RK4 step of size $h$ (Chapter~\ref{ch:integrators}) maps
$\vv{y}_q \mapsto \mat{P}\vv{y}_q$ with $\mat{P} = \sum_{k=0}^{4}
(h\mat{\Omega}/2)^k / k!$, the degree-4 truncation of the exponential ---
the stability polynomial $R(z)$ of RK4 evaluated at the matrix argument
(Hairer and Wanner~\cite{hairer1996}, ch.~IV). Using
equation~\eqref{eq:rigidbody:bigomega}, $\mat{P} = a\,\mat{I}_4 +
b\,\hat{\mat{\Omega}}$ with $\theta = h\norm{\vv{\omega}}/2$,
$\hat{\mat{\Omega}} = \mat{\Omega}/\norm{\vv{\omega}}$ (orthogonal and
skew), $a = 1 - \theta^2/2 + \theta^4/24$, and $b = \theta - \theta^3/6$.
Because $\vv{y}^{\mathsf T}\hat{\mat{\Omega}}\vv{y} = 0$ and
$\norm{\hat{\mat{\Omega}}\vv{y}} = \norm{\vv{y}}$, every step scales the
norm by exactly
\begin{equation}
  \rho(\theta) = \sqrt{a^2 + b^2}
    = \sqrt{1 - \frac{\theta^6}{72} + \frac{\theta^8}{576}}
    = 1 - \frac{\theta^6}{144} + \mathcal{O}(\theta^8),
  \label{eq:rigidbody:rk4norm}
\end{equation}
so after $N$ fixed steps the drift is, to leading order,
\begin{equation}
  \bigl|\norm{\quat{q}_N} - 1\bigr|
    = N\,\frac{\theta^6}{144}
      \bigl(1 + \mathcal{O}(\theta^2)\bigr),
  \qquad
  \theta = \frac{h \norm{\vv{\omega}}}{2} .
  \label{eq:rigidbody:normdrift}
\end{equation}
This bound is exact for constant $\vv{\omega}$ and is machine-checked by
\testid{rigidbody_quaternion_norm_drift}. For general motion under the
adaptive RKF7(8) driver, the per-step norm change is bounded by that
step's local error, which the controller holds at the group tolerances
(Chapter~\ref{ch:integrators}), so the drift accumulates at most linearly:
$|\norm{\quat{q}} - 1| \lesssim N_{\mathrm{steps}} (\mathrm{atol}_q +
\mathrm{rtol})$; the coning acceptance test measures it. Post-step
normalization $\quat{q} \mapsto \quat{q}/\norm{\quat{q}}$ removes the
drift entirely while leaving the represented attitude \emph{exactly}
unchanged (the DCM of equation~\eqref{eq:notation:quat2dcm} is degree-2
homogeneous, so rescaling touches only its overall $\norm{\quat{q}}^2$
factor), which is why FR-1 mandates it as the norm-control policy rather
than any projection that would alter the attitude.

\section{Discretization and implementation}
\label{sec:rigidbody:implementation}

There is no discretization: the module is a pointwise algebraic map,
evaluated inside the vehicle right-hand side at whatever states the
integrator (Chapter~\ref{ch:integrators}) requests. Implementation notes:
\begin{enumerate}
  \item \textbf{Module and traceability.} The model lives in
    \texttt{cpp/src/models/rigidbody.cpp}. The source comments echo the
    labels \texttt{eq:rigidbody:qdot}, \texttt{eq:rigidbody:euler},
    \texttt{eq:rigidbody:omegadot}, and \texttt{eq:rigidbody:normdrift}
    verbatim.
  \item \textbf{Packed state slice.} The raw-pointer entry point
    \texttt{rigidbody\_rhs} operates on the seven-component slice
    $[q_w, q_x, q_y, q_z, \omega_x, \omega_y, \omega_z]$ (scalar-first,
    per Table~\ref{tab:notation:eigenmap} rule 1), the layout the Phase~4
    vehicle state vector embeds.
  \item \textbf{Explicit $3 \times 3$ inverse.} The solve in
    equation~\eqref{eq:rigidbody:omegadot} uses the closed-form cofactor
    inverse (\texttt{Eigen::Matrix3d::inverse()}): IEEE basic operations
    only, in a fixed order, no pivot search, no allocation --- the
    relative error is bounded by a small multiple of
    $\kappa(\mat{I})\,\varepsilon$, and physical inertia tensors are
    well conditioned ($\kappa < 10^2$ for any credible vehicle).
  \item \textbf{No normalization inside the right-hand side.} Normalizing
    $\quat{q}$ inside the RHS would change the ODE the integrator sees and
    corrupt the embedded error estimate; by
    equation~\eqref{eq:rigidbody:qnorm} it is also unnecessary.
    Normalization is applied \emph{between} steps by the propagation
    layer via \texttt{rigidbody\_renormalize}, which returns the signed
    pre-normalization deviation $\norm{\quat{q}} - 1$ so callers can
    monitor the documented drift bound~\eqref{eq:rigidbody:normdrift}.
  \item \textbf{Determinism.} No heap allocation, single-threaded, fixed
    evaluation order, IEEE basic operations only (D-10); the golden
    generator's mpmath mirror reproduces the same formulation from the
    same committed binary64 inputs.
\end{enumerate}

\section{Validation evidence}
\label{sec:rigidbody:validation}

Table~\ref{tab:rigidbody:validation} lists the tests that constitute this
chapter's validation evidence. Each test identifier is machine-checked to
exist in the test tree by \texttt{scripts/lint\_chapters.py}; golden
provenance and tolerance derivations are recorded in
\texttt{tests/golden/attitude/manifest.toml}.

\begin{table}[htbp]
  \centering
  \caption{Validation evidence for the rigid-body attitude model (doctest,
    \texttt{cpp/tests/test\_rigidbody.cpp}).}
  \label{tab:rigidbody:validation}
  \begin{tabular}{lp{8.6cm}}
    \toprule
    Test ID & Acceptance basis \\
    \midrule
    \testid{rigidbody_rhs_golden} &
      Equations~\eqref{eq:rigidbody:qdot}
      and~\eqref{eq:rigidbody:omegadot} against extended-precision
      references (mpmath, 60 significant digits) at six committed states
      covering full inertia tensors, nonzero $\dot{\mat{I}}$, and torque:
      norm-relative $\le 10^{-12}$; structural exact-zero outputs
      bit-exact. \\
    \testid{rigidbody_zero_torque_properties} &
      Exact structural zeros: principal-axis spin and power-of-two
      spherical inertia give $\dot{\vv{\omega}} = \vv{0}$ exactly, a rest
      state gives $\dot{\quat{q}} = 0$ exactly; and under RK4 integration
      of the spherical case $\vv{\omega}$ stays bitwise constant while
      the attitude precesses. \\
    \testid{rigidbody_coning_closed_form} &
      Phase~4 exit criterion 4 (coning clause): RKF7(8) propagation of
      equations~\eqref{eq:rigidbody:qdot}--\eqref{eq:rigidbody:euler}
      against the closed form~\eqref{eq:rigidbody:coning} at eight
      checkpoints over five precession periods; attitude error angle
      $\le 10^{-9}$ rad and rate error $\le 10^{-9}$ norm-relative
      (measured worst case captured in the test output). \\
    \testid{rigidbody_intermediate_axis_flip} &
      Phase~4 exit criterion 4 (flip clause): across an intermediate-axis
      flip, $\vv{H}^{I}$ (magnitude and direction,
      equation~\eqref{eq:rigidbody:H}) and $T$
      (equation~\eqref{eq:rigidbody:T}) conserved to $10^{-10}$ relative
      at every accepted step; early-time rates match the independent
      mpmath trajectory to $10^{-9}$ norm-relative within the
      $e^{\lambda t}$ horizon of
      equation~\eqref{eq:rigidbody:lambda}; the flip itself reproduced
      ($\min \omega_2 < -0.9\,\omega_2(0)$). \\
    \testid{rigidbody_quaternion_norm_drift} &
      FR-1 norm-drift bound: measured RK4 drift matches
      equation~\eqref{eq:rigidbody:normdrift} to within 5\,\% over 400
      fixed steps of the constant-rate spin; adaptive-driver drift stays
      below the linear-accumulation bound; \texttt{rigidbody\_renormalize}
      restores unit norm to machine rounding and leaves the rate slice
      and the represented attitude untouched. \\
    \bottomrule
  \end{tabular}
\end{table}

\section{References}
\label{sec:rigidbody:references}

This chapter relies on Markley and Crassidis~\cite{markley2014} for the
standard forms of the quaternion kinematics, Euler's rotational equation,
and the classical torque-free axisymmetric solution (all re-derived above
in the project convention), and on Hairer and Wanner~\cite{hairer1996} for
the Runge--Kutta stability-function machinery behind the norm-drift bound,
with the local-error accumulation argument per Hairer, N{\o}rsett, and
Wanner~\cite{hairer1993}. Full bibliographic details appear in the
bibliography.
